% Kapitel 1 - Einleitung
\section{Introduction}
\label{key}
\onehalfspacing
{\large This chapter will give a small introduction on what the thesis and the implemented
	application is about. This covers an explanation for the application’s motivation and
	the environment the application will be used in.}


\subsection{Motivation}
{\large 
	
	The continuous monitoring of electrocardiographic (ECG) signals plays a crucial role in modern medicine, especially in the early detection and diagnosis of heart disease. While traditional ECG devices are often stationary or bulky, mobile technologies enable flexible and location-independent measurement. This opens up new possibilities for preventive health monitoring and telemedical care of patients.
	With the increasing use of wearable sensors and Bluetooth Low Energy (BLE) technology, new challenges arise in real-time data processing and visualization. A reliable and efficient software solution is required to capture, analyze and clearly display large amounts of ECG data in real time.
	The aim of this work is to develop a modern and high-performance front-end application that enables precise and intuitive visualization of ECG data.  This work thus makes an important contribution to digital healthcare and demonstrates the potential of mobile technologies for cardiology diagnostics.

	Furthermore the goal  is to  create a new version of an Android App once developed by fellow students from the faculty of Electrical and Information Technology. It's core functionallity is to receive and display ECG data in real time, via BLE. With the completion of this work, the new version will be able to receive and display ECG data in real time, via BLE. In addition it will be capable of an 24-hour mode which records ECG data for up to 24 hours, so the collected data can be analyzed later.
}
\subsection{Scope and Framework}
{\large
	This paper only deals with the app development process of the new version. Therefore, any implementation of the previous work done by the students is not going to be discussed. The hardware implementation is also not an element of this work. Generally speaking, this paper covers all topics that are essential for the development of the new app version. Specifically, technology stack of choice, IDE choice desicion, communication with hardware via Bluetooth low energy, real-time data monitoring, implementing the 24 hour mode.  In addition, the outlook explains procedures for the implementation of future features that could not be realized within the limits of this paper.
		
}
\subsection{Procedure}
{\large
	The second chapter, \textit{Technical Fundamentals}, explores the Android architecture, with a particular focus on the Java Framework API. 
	Giving valuable insights into the core functionality of the operating system aswell as an glimpse into the BLE-Stack, laying the groundwork for a better understanding of the implementation details presented later.  
	Chapter three delves into the functional requirements, systematically breaking them down into distinct sub-requirements. These requirements are revisited in chapter four, where they are translated into actual code, accompanied by relevant code snippets to illustrate their implementation.
}
